\documentclass[11pt]{article}
\usepackage{amsmath, amsthm, amssymb}
\usepackage{mathtools}
\usepackage[margin=1in]{geometry}
\usepackage{enumitem}
\usepackage{hyperref}

\theoremstyle{plain}
\newtheorem{theorem}{Theorem}
\newtheorem{lemma}[theorem]{Lemma}
\newtheorem{proposition}[theorem]{Proposition}
\newtheorem{corollary}[theorem]{Corollary}

\theoremstyle{definition}
\newtheorem{definition}[theorem]{Definition}
\newtheorem{example}[theorem]{Example}

\theoremstyle{remark}
\newtheorem{remark}[theorem]{Remark}
\newtheorem{question}[theorem]{Question}

\newcommand{\Z}{\mathbb{Z}}
\newcommand{\N}{\mathbb{N}}
\newcommand{\Q}{\mathbb{Q}}
\newcommand{\Sym}{\mathfrak{S}}
\newcommand{\Pibs}{\Pi_q}
\newcommand{\Vlam}{V_\lambda}
\newcommand{\trace}{\mathrm{tr}}
\newcommand{\PhiThree}{\Phi_3}

\title{Twin identity and $3q\Phi_3$-perturbation rule \\
       for $\sigma_1$-G1 cores at $n = 6$}
\author{Clio}
\date{29 April 2026}

\begin{document}
\maketitle

\begin{abstract}
Let $A_\lambda(q) \in \Z_{\geq 0}[q]$ denote the palindromic core of the staircase
Bott--Samelson trace $\sigma_1(\Vlam) = \trace(\Pibs|_\Vlam)$ on the Specht module
$\Vlam$, after extracting $q^{a_\lambda}(1+q)^{b_\lambda}$.  We prove three
polynomial identities about the cores at $n = 5$ and $n = 6$, organised around
the key shape $\lambda = (3,1,1)$ whose core $A_{(3,1,1)}(q) = q^2 + 4q + 1$ acts
as a multiplicative scaffold:
\begin{align*}
\textsc{Lemma (n=5, atom relations)} & : \quad
  A_{(3,2)} = A_{(3,1,1)}^2 + q\PhiThree, \qquad
  A_{(4,1)} = A_{(3,1,1)}^2 - q\PhiThree, \\[4pt]
\textsc{Theorem (n=6, twin pair)} & : \quad
  A_{(3,2,1)} = A_{(3,1,1)}^3 + 3q\PhiThree \cdot A_{(3,1,1)}, \\
  & \phantom{ : \quad} A_{(5,1)} = A_{(3,1,1)}^3 - 3q\PhiThree \cdot A_{(3,1,1)},
\end{align*}
where $\PhiThree(q) = q^2 + q + 1$ is the third cyclotomic polynomial.  The three
target identities of the proof session (Theorems A, B, C of \texttt{PROVE.md})
follow as immediate corollaries.  In particular,
\[
A_{(3,2,1)} + A_{(5,1)} = 2 A_{(3,1,1)}^3,
\qquad
A_{(3,2,1)} - A_{(5,1)} = 6q\PhiThree \cdot A_{(3,1,1)}.
\]
The proof of each identity is a direct polynomial verification using the
stair-basis $\sigma_1$-G1 expansion of \cite{Apr24}.  The structural
\emph{origin} of the symmetric $\pm q\PhiThree$ pattern at $n=5$, and of its
exact triplication at $n=6$, remains mysterious; we close with the precise
combinatorial questions left open.
\end{abstract}

\section{Setup}

\subsection{Hecke algebra and $\sigma_1$}

Let $H_q$ be the Hecke algebra of $\Sym_n$ over $\Z[q,q^{-1}]$, with generators
$T_1, \ldots, T_{n-1}$ subject to $T_i^2 = (q-1)T_i + q$.  Fix the staircase
reduced word for the longest element $w_0 \in \Sym_n$:
\[
\underline{w}_0 = (i_1, \ldots, i_N), \qquad N = \binom{n}{2},
\]
and define the staircase Bott--Samelson product
\[
\Pibs := \prod_{j=1}^N (T_{i_j} + 1) \in H_q.
\]
For each partition $\lambda \vdash n$ let $\Vlam$ be the (cell/Specht) module and
\[
\sigma_1(\Vlam) := \trace(\Pibs|_\Vlam) \in \Z[q].
\]
By the bar involution $\sigma_1$ is palindromic of degree $\leq N$ in $q$, and by
Kazhdan--Lusztig positivity $\sigma_1 \in \Z_{\geq 0}[q]$.  Let $a_\lambda$,
$b_\lambda$ be the maximal $q$- and $(1+q)$-divisibilities of $\sigma_1(\Vlam)$,
and define the \emph{$\sigma_1$-G1 core}
\[
A_\lambda(q) := \frac{\sigma_1(\Vlam)}{q^{a_\lambda}(1+q)^{b_\lambda}} \in \Z[q].
\]

\subsection{The $\sigma_1$-G1 stair basis}

\begin{theorem}[\cite{Apr24}, $\sigma_1$-G1]\label{thm:G1}
For every $\lambda \vdash n$ there exist non-negative integers $n_k(\lambda)
\in \Z_{\geq 0}$ for $k \in \{0,1,\ldots,N\}$ with $k \equiv N \pmod 2$, such that
\begin{equation}
\sigma_1(\Vlam) = \sum_{\substack{0 \leq k \leq N \\ k \equiv N \!\!\!\!\pmod 2}}
n_k(\lambda)\,(1+q)^k\,q^{(N-k)/2}.
\label{eq:G1}
\end{equation}
The numbers $n_k(\lambda)$ are obtained as a sum over closed length-$N$ paths in
the Kazhdan--Lusztig $W$-graph of $\Vlam$, with $k$ ``stay'' (descent) steps and
$N-k$ ``move'' (KL-edge) steps.
\end{theorem}

Dividing~\eqref{eq:G1} by $q^{a_\lambda}(1+q)^{b_\lambda}$ gives the
\emph{stair-basis expansion} of $A_\lambda$:
\begin{equation}
A_\lambda(q) = \sum_{i=0}^{r_\lambda} c_i(\lambda)\,
   q^{r_\lambda-i}\,(1+q)^{2i},
\qquad c_i(\lambda) := n_{b_\lambda + 2i}(\lambda),
\label{eq:stairA}
\end{equation}
where $r_\lambda$ is the integer such that the support of $n_k(\lambda)$ is
$\{b_\lambda, b_\lambda + 2, \ldots, b_\lambda + 2 r_\lambda\}$.  Equivalently
$r_\lambda = (N - 2 a_\lambda - b_\lambda)/2$.  All $c_i(\lambda) \in
\Z_{\geq 0}$ by Theorem~\ref{thm:G1}.  Two examples used throughout:

\begin{example}\label{ex:n5basis}
At $n = 5$ ($N = 10$): each rank-$1$ atom has $r_\lambda = 2$, so the stair
basis at $n=5$ is $(q^2,\, q(1+q)^2,\, (1+q)^4)$.  We have
\[
A_{(3,2)} = 3\,q^2 + 5\,q(1+q)^2 + (1+q)^4, \qquad
A_{(4,1)} = 5\,q^2 + 3\,q(1+q)^2 + (1+q)^4,
\]
\[
A_{(3,1,1)} = 2\,q + (1+q)^2 \quad \text{(at $r_\lambda = 1$, basis $(q,(1+q)^2)$).}
\]
\end{example}

\begin{example}\label{ex:n6basis}
At $n = 6$ ($N = 15$): the rank-$2$ atoms $A_{(3,2,1)}$ and $A_{(5,1)}$ both
have $r_\lambda = 3$, so the stair basis at $n=6$ for them is
$(q^3,\, q^2(1+q)^2,\, q(1+q)^4,\, (1+q)^6)$.  We have
\[
A_{(3,2,1)} = 2\,q^3 + 15\,q^2(1+q)^2 + 9\,q(1+q)^4 + (1+q)^6,
\]
\[
A_{(5,1)} = 14\,q^3 + 9\,q^2(1+q)^2 + 3\,q(1+q)^4 + (1+q)^6.
\]
\end{example}

\subsection{Five base atoms}

The polynomials at the centre of this paper are:
\begin{align*}
A_{(3,1,1)}  &= q^2 + 4q + 1, \\
A_{(3,2)}    &= q^4 + 9q^3 + 19q^2 + 9q + 1, \\
A_{(4,1)}    &= q^4 + 7q^3 + 17q^2 + 7q + 1, \\
A_{(3,2,1)}  &= q^6 + 15q^5 + 66q^4 + 106q^3 + 66q^2 + 15q + 1, \\
A_{(5,1)}    &= q^6 + 9q^5  + 36q^4 + 70q^3  + 36q^2 + 9q  + 1.
\end{align*}
We write $\PhiThree(q) := q^2 + q + 1$ for the third cyclotomic polynomial.

\section{The $\pm q\PhiThree$ pattern at $n = 5$}

\begin{lemma}[$n=5$ atom relations]\label{lem:n5}
\[
A_{(3,2)}(q)  = A_{(3,1,1)}(q)^2 + q\,\PhiThree(q), \qquad
A_{(4,1)}(q)  = A_{(3,1,1)}(q)^2 - q\,\PhiThree(q).
\]
\end{lemma}

\begin{proof}
Direct computation:
\begin{align*}
A_{(3,1,1)}^2 &= (q^2+4q+1)^2 = q^4 + 8q^3 + 18q^2 + 8q + 1, \\
q\,\PhiThree(q) &= q\,(q^2+q+1) = q^3 + q^2 + q.
\end{align*}
Adding: $A_{(3,1,1)}^2 + q\,\PhiThree = q^4 + 9q^3 + 19q^2 + 9q + 1 = A_{(3,2)}$.
Subtracting: $A_{(3,1,1)}^2 - q\,\PhiThree = q^4 + 7q^3 + 17q^2 + 7q + 1 =
A_{(4,1)}$.\qedhere
\end{proof}

\begin{corollary}[$n = 5$ symmetric averaging]\label{cor:n5sum}
\[
A_{(3,2)} + A_{(4,1)} = 2\,A_{(3,1,1)}^2,
\qquad
A_{(3,2)} - A_{(4,1)} = 2\,q\,\PhiThree(q) =: \delta(q).
\]
\end{corollary}

\begin{remark}[Stair-basis interpretation of $q\PhiThree$]
\label{rem:stairperturb}
In the $n=5$ stair basis $(q^2, q(1+q)^2, (1+q)^4)$, $A_{(3,1,1)}^2$ has
coordinates $(4,4,1)$, while $A_{(3,2)}$ has $(3,5,1)$ and $A_{(4,1)}$ has
$(5,3,1)$.  The perturbation $q\PhiThree = -q^2 + q(1+q)^2$ has stair
coordinates $(-1, +1, 0)$, so $A_{(3,2)} - A_{(3,1,1)}^2$ flips $-1$ from $q^2$
into $+1$ for $q(1+q)^2$, and $A_{(4,1)} - A_{(3,1,1)}^2$ does the reverse.
Lemma~\ref{lem:n5} is the statement that the rank-$1$ pair
$\{(3,2), (4,1)\}$ at $n = 5$ is symmetric across $A_{(3,1,1)}^2$ along this
$q\PhiThree$-axis.
\end{remark}

\section{The $\pm 3q\PhiThree$ pattern at $n = 6$ (twin pair)}

\begin{theorem}[$n=6$ twin pair]\label{thm:n6twin}
\[
A_{(3,2,1)}(q) = A_{(3,1,1)}(q)^3 + 3\,q\,\PhiThree(q)\cdot A_{(3,1,1)}(q),
\]
\[
A_{(5,1)}(q)   = A_{(3,1,1)}(q)^3 - 3\,q\,\PhiThree(q)\cdot A_{(3,1,1)}(q).
\]
Equivalently,
\[
\frac{A_{(3,2,1)}(q)}{A_{(3,1,1)}(q)} = A_{(3,1,1)}(q)^2 + 3\,q\,\PhiThree(q),
\qquad
\frac{A_{(5,1)}(q)}{A_{(3,1,1)}(q)} = A_{(3,1,1)}(q)^2 - 3\,q\,\PhiThree(q).
\]
\end{theorem}

\begin{proof}
Compute
\[
A_{(3,1,1)}^3 = (q^2+4q+1)^3 = q^6 + 12q^5 + 51q^4 + 88q^3 + 51q^2 + 12q + 1.
\]
And $3 q \PhiThree \cdot A_{(3,1,1)} = 3q(q^2+q+1)(q^2+4q+1) = 3 q
(q^4 + 5q^3 + 7q^2 + 5q + 1) = 3q^5 + 15q^4 + 21q^3 + 15q^2 + 3q$.
Adding,
\[
A_{(3,1,1)}^3 + 3 q \PhiThree A_{(3,1,1)}
= q^6 + 15q^5 + 66q^4 + 106q^3 + 66q^2 + 15q + 1
= A_{(3,2,1)}.
\]
Subtracting,
\[
A_{(3,1,1)}^3 - 3 q \PhiThree A_{(3,1,1)}
= q^6 + 9q^5 + 36q^4 + 70q^3 + 36q^2 + 9q + 1
= A_{(5,1)}.
\qedhere
\]
\end{proof}

\begin{corollary}[Twin sum and difference]\label{cor:twinsum}
\[
A_{(3,2,1)} + A_{(5,1)} = 2\,A_{(3,1,1)}^3,
\qquad
A_{(3,2,1)} - A_{(5,1)} = 6\,q\,\PhiThree \cdot A_{(3,1,1)} = 3\,\delta \cdot A_{(3,1,1)}.
\]
\end{corollary}

\begin{remark}[Stair-basis interpretation]
\label{rem:stairperturb6}
In the $n=6$ stair basis $(q^3, q^2(1+q)^2, q(1+q)^4, (1+q)^6)$,
$A_{(3,1,1)}^3$ has stair coordinates $(8, 12, 6, 1)$, while $A_{(3,2,1)}$ has
$(2, 15, 9, 1)$ and $A_{(5,1)}$ has $(14, 9, 3, 1)$.  The perturbation
$3q\PhiThree \cdot A_{(3,1,1)}$ has stair coordinates $(-6, +3, +3, 0)$, the
Cauchy product of the $n=5$ perturbation $3q\PhiThree = (-3, +3, 0)$ and
$A_{(3,1,1)} = (2, 1)$.  The pair $\{(3,2,1), (5,1)\}$ at $n=6$ is symmetric
across $A_{(3,1,1)}^3$ along an enlarged perturbation axis.
\end{remark}

\section{Theorems A, B, C of \texttt{PROVE.md} as corollaries}

\begin{corollary}[Theorem A, Twin identity]\label{cor:A}
\[
A_{(3,2,1)} + A_{(5,1)} = A_{(3,1,1)}\cdot \bigl(A_{(3,2)} + A_{(4,1)}\bigr).
\]
\end{corollary}

\begin{proof}
By Corollary~\ref{cor:twinsum}, $A_{(3,2,1)} + A_{(5,1)} = 2 A_{(3,1,1)}^3
= A_{(3,1,1)}\cdot 2 A_{(3,1,1)}^2$.  By Corollary~\ref{cor:n5sum},
$2 A_{(3,1,1)}^2 = A_{(3,2)} + A_{(4,1)}$.\qedhere
\end{proof}

\begin{corollary}[Theorem B, Sign-flipped pair]\label{cor:B}
With $\delta(q) = A_{(3,2)} - A_{(4,1)} = 2q\PhiThree(q)$,
\[
A_{(3,2,1)} = A_{(3,1,1)}\cdot(A_{(3,2)} + \delta) = A_{(3,1,1)}\cdot(2 A_{(3,2)} - A_{(4,1)}),
\]
\[
A_{(5,1)}  = A_{(3,1,1)}\cdot(A_{(4,1)} - \delta) = A_{(3,1,1)}\cdot(2 A_{(4,1)} - A_{(3,2)}).
\]
\end{corollary}

\begin{proof}
By Lemma~\ref{lem:n5}, $A_{(3,2)} = A_{(3,1,1)}^2 + q\PhiThree$, so
$A_{(3,2)} + \delta = A_{(3,1,1)}^2 + q\PhiThree + 2 q\PhiThree =
A_{(3,1,1)}^2 + 3 q\PhiThree$.  Multiplying by $A_{(3,1,1)}$ and using
Theorem~\ref{thm:n6twin} gives $A_{(3,1,1)}\cdot (A_{(3,2)} + \delta) =
A_{(3,1,1)}^3 + 3q\PhiThree A_{(3,1,1)} = A_{(3,2,1)}$.  The substitution
$A_{(3,2)} + \delta = 2 A_{(3,2)} - A_{(4,1)}$ follows from
$\delta = A_{(3,2)} - A_{(4,1)}$ (Corollary~\ref{cor:n5sum}).  Symmetric
argument for $A_{(5,1)}$.\qedhere
\end{proof}

\begin{corollary}[Theorem C, $2q$-perturbation rule]\label{cor:C}
$A_{(3,2,1)} = A_{(3,1,1)}\cdot \mathrm{atom}_{\mathrm{RTL}}$, where
$\mathrm{atom}_{\mathrm{RTL}}(q) := q^4 + 11 q^3 + 21 q^2 + 11 q + 1$ satisfies
\[
\mathrm{atom}_{\mathrm{RTL}}(q) = A_{(3,2)}(q) + 2\,q\,\PhiThree(q)
                                = A_{(3,1,1)}(q)^2 + 3\,q\,\PhiThree(q).
\]
\end{corollary}

\begin{proof}
The left equality is Corollary~\ref{cor:B} (substitute $A_{(3,2)} + \delta$).
For the right equality, apply Lemma~\ref{lem:n5} again:
$A_{(3,2)} + 2q\PhiThree = A_{(3,1,1)}^2 + q\PhiThree + 2q\PhiThree = A_{(3,1,1)}^2 + 3q\PhiThree$.\qedhere
\end{proof}

\section{Computational verification}

All polynomial identities above were verified by direct expansion in
\texttt{sympy}; see
\path{~/projects/scratch/2026-04-29-twin-identity-proof/verify.py}.  Specifically:
\begin{itemize}[leftmargin=2em]
\item Lemma~\ref{lem:n5} (two identities): \textbf{OK}.
\item Corollary~\ref{cor:n5sum} (sum, difference): \textbf{OK}.
\item Theorem~\ref{thm:n6twin} (two identities): \textbf{OK}.
\item Corollary~\ref{cor:twinsum} (sum = $2 A_{(3,1,1)}^3$, difference =
      $6q\PhiThree A_{(3,1,1)}$): \textbf{OK}.
\item Corollary~\ref{cor:A} (Theorem A): \textbf{OK}.
\item Corollary~\ref{cor:B} (Theorem B, both halves): \textbf{OK}.
\item Corollary~\ref{cor:C} (Theorem C, both equalities): \textbf{OK}.
\end{itemize}
The five base atoms have themselves been computed independently from the
Hoefsmit/seminormal Hecke matrices (see
\path{~/projects/scratch/2026-04-25-hu-mathas-test/compute_filtration.py}
and \path{~/projects/scratch/2026-04-29-branching-decomp-other-shapes/compute.py}).

\section{What is structural, what is mysterious}

\subsection{What is structural}

\paragraph{Stair-basis expansion.}
The $\sigma_1$-G1 theorem (\cite{Apr24}, recalled here as
Theorem~\ref{thm:G1}) gives every $A_\lambda$ as a non-negative integer
combination of stair monomials $q^{r-i}(1+q)^{2i}$, $0 \leq i \leq r$.  This is a
representation-theoretic statement: the integer $n_k(\lambda)$ counts closed
length-$N$ paths in the $W$-graph of $\Vlam$ with exactly $k$ ``stay''
(descent) steps and $N-k$ ``move'' (KL-edge) steps; the coefficient of each
path is a product of KL $\mu$-coefficients and is non-negative by KL positivity.

\paragraph{Multiplicativity of stair basis.}
The stair monomials of degree $r$ at $n$ multiply with the stair monomials of
degree $r'$ at $n'$ as $q^{r-i}(1+q)^{2i} \cdot q^{r'-j}(1+q)^{2j} =
q^{(r+r') - (i+j)}(1+q)^{2(i+j)}$, i.e., the stair basis is closed under
polynomial multiplication.  This is the only structural fact used to lift
identities from $n=5$ atoms to $n=6$ atoms: the polynomial product
$A_{(3,1,1)} \cdot A_{(3,2)}$ is again a non-negative integer combination of
stair monomials, in the stair basis at degree $r_{(3,1,1)} + r_{(3,2)}$.

\paragraph{Symmetric pair across a power of $A_{(3,1,1)}$.}
The empirical pattern is that the symmetric pair $\{(3,2), (4,1)\}$ at $n=5$
straddles $A_{(3,1,1)}^2$, and the symmetric pair $\{(3,2,1), (5,1)\}$ at $n=6$
straddles $A_{(3,1,1)}^3$.  Each is a $\pm$ perturbation by an integer multiple
of $q\PhiThree(q)$ times an appropriate power of $A_{(3,1,1)}$.

\subsection{What remains mysterious}

\begin{question}[Origin of $q\PhiThree$]
Why is the perturbation kernel $q\PhiThree(q) = q + q^2 + q^3$?  Both
$\PhiThree$ and the leading $q$ are natural in the $W$-graph context
($\PhiThree$ specialises the Hecke algebra at a primitive third root of unity,
where new idempotents appear; $q$ is one full ``move'' weight in the
$\sigma_1$-G1 expansion).  But we have not identified the specific closed
$W$-graph paths whose generating function is $q\PhiThree$.  In particular, the
$n=5$ stair-basis decomposition $q\PhiThree = -q^2 + q(1+q)^2 + 0\cdot(1+q)^4$
involves a \emph{negative} stair coefficient, which is necessarily a
\emph{signed} cancellation between paths in the $W$-graphs of $V_{(3,2)}$ and
$V_{(4,1)}$, not a path-counting identity in either separately.
\end{question}

\begin{question}[Origin of the factor $3$]
Why does the $n=5$ kernel $q\PhiThree$ \emph{exactly triple} on passing from
$n=5$ to the lifted $n=6$ pair, becoming $3 q\PhiThree$ when measured against
$A_{(3,1,1)}$ as multiplicative scaffold?
\begin{itemize}[leftmargin=2em]
    \item The number of removable corners of $(3,2,1)$ at $n=6$ is $3$ (cells
          $(1,3), (2,2), (3,1)$); this matches.  But the same factor $3$ also
          appears in $A_{(5,1)}/A_{(3,1,1)} = A_{(3,1,1)}^2 - 3q\PhiThree$,
          and $(5,1)$ has only $2$ removable corners.  So ``corners'' is not
          the right invariant.
    \item Equivalently, in the stair basis, the kernel passes from $(-1, +1, 0)$
          at $n=5$ to $(-3, +3, 0)$ at $n=5$-rescaled, and then is convolved with
          $A_{(3,1,1)} = (2, 1)$ to produce the $n=6$ shift $(-6, +3, +3, 0)$.
          What is the $n=6$ structural object that scales the kernel by $3$?
\end{itemize}
\end{question}

\begin{question}[Reason $A_{(3,1,1)}$ divides $A_{(5,1)}$]
The divisibility $A_{(3,1,1)} \mid A_{(3,2,1)}$ is consistent with branching:
$V_{(3,2,1)}\!\downarrow\!\Sym_5 \supset V_{(3,1,1)}$.  But
$V_{(5,1)}\!\downarrow\!\Sym_5 = V_{(5)} \oplus V_{(4,1)}$ does \emph{not}
contain $V_{(3,1,1)}$, yet $A_{(3,1,1)}$ still divides $A_{(5,1)}$.  Hence
$A_{(3,1,1)}$ plays a structural role at $n=6$ that transcends naive Hecke
branching.  We have no explanation for this beyond the polynomial identity.
\end{question}

\begin{remark}[Specialisation at $q = \omega$, primitive third root of unity]
\label{rem:omega}
Since $\PhiThree(\omega) = 0$, both $q\PhiThree$ and $3 q \PhiThree$ vanish at
$q = \omega$.  Therefore at $q = \omega$,
\[
A_{(3,2)}(\omega) = A_{(4,1)}(\omega) = A_{(3,1,1)}(\omega)^2,
\qquad
A_{(3,2,1)}(\omega) = A_{(5,1)}(\omega) = A_{(3,1,1)}(\omega)^3.
\]
A direct computation gives $A_{(3,1,1)}(\omega) = 3 \omega$, so
$A_{(3,2,1)}(\omega) = A_{(5,1)}(\omega) = 27 \omega^3 = 27$.  This collapse at
$q = \omega$ is consistent with the perturbation pattern but, again, not
explained by it.
\end{remark}

\subsection{What is the gap}

A \emph{structural} proof of Lemma~\ref{lem:n5} and Theorem~\ref{thm:n6twin}
would identify, for each of the four directed perturbation polynomials
\[
A_{(3,2)} - A_{(3,1,1)}^2 = +q\PhiThree, \;\;\;
A_{(4,1)} - A_{(3,1,1)}^2 = -q\PhiThree,
\]
\[
A_{(3,2,1)}/A_{(3,1,1)} - A_{(3,1,1)}^2 = +3q\PhiThree, \;\;\;
A_{(5,1)}/A_{(3,1,1)} - A_{(3,1,1)}^2 = -3q\PhiThree,
\]
a precise object on which the difference is computed: a class of $W$-graph
paths, a Soergel bimodule difference, an $\mathrm{Ext}$-group, or a similar
categorical or representation-theoretic invariant whose generating function is
$\pm q\PhiThree$ at $n=5$ and $\pm 3 q\PhiThree A_{(3,1,1)}$ at $n=6$.  We do
not produce such an object here.  The polynomial proofs above are honest
identities, but they are not yet proofs at the $W$-graph or Hecke level.

\section{Discussion}

\subsection{Relation to other rank-$2$ shapes at $n=6$}

The third rank-$2$ shape at $n=6$ is~$(4,1,1)$ (whose conjugate is
$(3,1,1,1)$); we have not addressed it here because, computationally,
$A_{(4,1,1)} = A_{(3,1,1)}\cdot A_{(3,1)\downarrow n=4}\cdot \cdots$ admits a
different multiplicative structure.  Likewise the rank-$3$ shape~$(4,2)$ does
\emph{not} factor through $A_{(3,1,1)}$ multiplicatively (verified
computationally: $A_{(3,1,1)}\nmid A_{(4,2)}$).  The dichotomy ``rank-$2$
factors through $A_{(3,1,1)}$, rank-$3$ does not'' was identified earlier (Apr
26 work, $\lambda=(4,2)$ at $n=6$).

\subsection{What this proof unlocks}

Theorem~\ref{thm:n6twin}, combined with the trace-square identity at
$\lambda = (3,2,1)$ (Apr 29 work,
\path{~/projects/scratch/2026-04-29-trace-square-321/}),
\[
A_{(3,2,1)}^2 = G_{(3,2,1)} + 2 q \cdot Q_{(3,2,1)}
\quad \text{where}\;\;
G_{(3,2,1)} = \sigma_2(\Vlam)/q^8(1+q)^2,
\]
and $Q_{(3,2,1)}$ is the deg-$10$ orphan atom $A_{10}(q)$, specialises to
\[
A_{(3,1,1)}^2\,(\mathrm{atom}_{\mathrm{RTL}})^2 = G_{(3,2,1)} + 2 q\,A_{10}.
\]
Substituting $\mathrm{atom}_{\mathrm{RTL}} = A_{(3,1,1)}^2 + 3 q\PhiThree$
(Corollary~\ref{cor:C}):
\[
A_{(3,1,1)}^2\,\bigl(A_{(3,1,1)}^2 + 3q\PhiThree\bigr)^2 = G_{(3,2,1)} + 2 q A_{10}.
\]
This is a polynomial identity now expressed entirely in $n=5$ atoms and the
universal kernel $q\PhiThree$.  It places $A_{10}$ as the obstruction to
$\sigma_2$-G1 inheritance at $\lambda = (3,2,1)$ admitting a closed-form
expression in $n=5$ atoms.  We do not pursue this further here.

\subsection{Stretch goal: generalisation to rank-$2$ pairs at all $n$}

The $(3,2)/(4,1)$ pair at $n=5$ and the $(3,2,1)/(5,1)$ pair at $n=6$ both
satisfy
\[
A_\lambda + A_{\lambda^*} = 2 A_{(3,1,1)}^k
\]
for $k = 2, 3$ respectively, where $A_{(3,1,1)} = q^2 + 4q + 1$.  The
conjecture
\[
A_\lambda + A_{\lambda^*} = 2 A_{(3,1,1)}^{k_n}
\]
for $k_n$ depending on $n$ would generalise to an infinite family.  Numerical
evidence at $n = 7$ for the candidate pair $(4,2,1) \leftrightarrow (5,1,1)$ would
falsify or confirm this.  We have not tested this here.

\begin{thebibliography}{99}
\bibitem{Apr24}
Clio, \emph{Non-negativity of the trace core for the staircase Bott--Samelson
operator on Hecke cell modules},
\path{~/projects/proofs/2026-04-24-sigma1-G1.tex},
24 April 2026.
\bibitem{KL}
D.~Kazhdan and G.~Lusztig, \emph{Representations of Coxeter groups and Hecke
algebras}, Invent. Math. 53 (1979), 165--184.
\bibitem{EW}
B.~Elias and G.~Williamson, \emph{The Hodge theory of Soergel bimodules}, Ann.
of Math. (2) 180 (2014), 1089--1136.
\end{thebibliography}

\end{document}
